\chapter{پردازش رشته‌ها و محاسبات عددی}

ماهیت برنامه‌های کامپیوتری اساساً بر تعامل با داده‌ها استوار است. در فصول پیشین، تمرکز ما بر پردازش داده‌ها در سطح فایل بود؛ با این حال، بسیاری از چالش‌های برنامه‌نویسی نیازمند کار با واحدهای کوچک‌تری از داده، یعنی رشته‌ها (\en{Strings}) و اعداد (\en{Numbers}) هستند.

در این فصل، به بررسی آن دسته از قابلیت‌های پوسته می‌پردازیم که برای پردازش رشته‌ها و محاسبات عددی طراحی شده‌اند. پوسته چندین روش برای بسط پارامتر (\en{Parameter Expansion}) ارائه می‌دهد که انجام عملیات روی رشته‌ها را تسهیل می‌کند. افزون بر بسط محاسباتی (\en{Arithmetic Expansion}) که در فصل ۷ به آن اشاره شد، ابزار خط فرمانی مشهوری به نام \cmd{bc} نیز معرفی می‌شود که برای محاسبات ریاضی در سطوح پیشرفته‌تر به کار می‌رود.

\section{بسط پارامتر (Parameter Expansion)}

اگرچه مفهوم بسط پارامتر در فصل ۷ مطرح شد، اما به دلیل کاربرد گسترده‌تر آن در اسکریپت‌نویسی نسبت به محیط مستقیم خط فرمان، به تفصیل بررسی نشد. پیش از این با اشکال ساده‌ای از بسط پارامتر، مانند فراخوانی متغیرهای پوسته، کار کرده‌ایم؛ اما پوسته امکانات بسیار وسیع‌تری در این زمینه فراهم می‌کند.

\begin{note}
\textbf{نکته:} همواره توصیه می‌شود بسط‌های پارامتر را درون دابل‌کوتیشن (\en{Double Quotes}) قرار دهید تا از بروز پدیده ناخواسته «تفکیک کلمات» (\en{Word Splitting}) جلوگیری شود؛ مگر آنکه دلیل فنی خاصی برای عدم استفاده از آن داشته باشید. این موضوع به‌ویژه هنگام تعامل با نام فایل‌هایی که حاوی فاصله (\en{Space}) یا کاراکترهای خاص هستند، اهمیت دوچندان می‌یابد.
\end{note}

\subsection{پارامترهای پایه}

ساده‌ترین شکل بسط پارامتر، همان فراخوانی معمول متغیرهاست؛ به عنوان مثال:

\begin{latin}
\begin{lstlisting}
$a
\end{lstlisting}
\end{latin}

هنگامی که این عبارت بسط داده می‌شود، پوسته آن را با مقدار ذخیره‌شده در متغیر \cmd{a} جایگزین می‌کند. پارامترهای ساده را می‌توان جهت وضوح بیشتر درون آکولاد نیز محصور کرد:

\begin{latin}
\begin{lstlisting}
${a}
\end{lstlisting}
\end{latin}

این نحو نگارش به تنهایی تغییری در نتیجه بسط ایجاد نمی‌کند، اما زمانی که متغیر بلافاصله در مجاورت نویسه‌های دیگر قرار می‌گیرد، استفاده از آکولاد برای جلوگیری از تشخیص اشتباه نام متغیر توسط پوسته الزامی است. به مثال زیر توجه کنید که در آن قصد داریم رشته \file{\_} را به انتهای محتوای متغیر \cmd{a} پیوند بزنیم:

\begin{latin}
\begin{lstlisting}
[me@linuxbox ~]$ a="foo"
[me@linuxbox ~]$ echo "$a_file"

\end{lstlisting}
\end{latin}

در صورت اجرای دستور فوق، هیچ خروجی مشاهده نخواهید کرد؛ زیرا پوسته به دنبال متغیری به نام \cmd{a\_file} می‌گردد که تعریف نشده است. برای حل این مشکل و تفکیک نام متغیر از متن مجاور، باید از آکولاد استفاده کرد:

\begin{latin}
\begin{lstlisting}
[me@linuxbox ~]$ echo "${a}_file"
foo_file
\end{lstlisting}
\end{latin}

علاوه بر این، پیش‌تر اشاره شد که برای دسترسی به پارامترهای موقعیتی (\en{Positional Parameters}) دو رقمی (بزرگ‌تر از ۹)، استفاده از آکولاد اجباری است. به عنوان مثال، برای فراخوانی یازدهمین آرگومان ورودی، باید به این صورت عمل کرد:

\begin{latin}
\begin{lstlisting}
${11}
\end{lstlisting}
\end{latin}

\subsection{مدیریت متغیرهای تعریف‌نشده یا تهی}

تعدادی از الگوهای بسط پارامتر به‌طور اختصاصی برای مواجهه با متغیرهای «تعریف‌نشده» (\en{Unset}) یا «تهی» (\en{Empty}) طراحی شده‌اند. این قابلیت‌ها جهت مدیریت پارامترهای موقعیتیِ ارائه‌نشده و تخصیص مقادیر پیش‌فرض (\en{Default Values}) به پارامترها بسیار کارآمد هستند.

\textbf{\cmd{\$\{parameter:-word\}}}

چنانچه \file{parameter} تعریف نشده یا مقدار آن تهی باشد، خروجی این بسط رشته \file{word} خواهد بود. در غیر این صورت، مقدار خودِ پارامتر بازگردانده می‌شود. به مثال زیر توجه کنید:

\begin{latin}
\begin{lstlisting}
[me@linuxbox ~]$ foo=
[me@linuxbox ~]$ echo ${foo:-"substitute value if unset"}
substitute value if unset
[me@linuxbox ~]$ echo $foo

[me@linuxbox ~]$ foo=bar
[me@linuxbox ~]$ echo ${foo:-"substitute value if unset"}
bar
[me@linuxbox ~]$ echo $foo
bar
\end{lstlisting}
\end{latin}

\textbf{\cmd{\$\{parameter:=word\}}}

این ساختار مشابه الگوی قبلی است، با این تفاوت که اگر \file{parameter} تعریف‌نشده یا تهی باشد، علاوه بر بازگرداندن مقدار \file{word}، این مقدار را به‌طور دائمی به خودِ پارامتر نیز اختصاص می‌دهد. در واقع، این روش برای مقداردهی اولیه متغیرهای فاقد مقدار استفاده می‌شود. اگر پارامتر از قبل دارای مقدار باشد، بسط با همان مقدار موجود خاتمه می‌یابد و تغییری در آن ایجاد نمی‌شود.

\begin{latin}
\begin{lstlisting}
[me@linuxbox ~]$ foo=
[me@linuxbox ~]$ echo ${foo:="default value if unset"}
default value if unset
[me@linuxbox ~]$ echo $foo
default value if unset
[me@linuxbox ~]$ foo=bar
[me@linuxbox ~]$ echo ${foo:="default value if unset"}
bar
[me@linuxbox ~]$ echo $foo
bar
\end{lstlisting}
\end{latin}

\begin{note}
\textbf{نکته:} توجه داشته باشید که پارامترهای موقعیتی (\en{Positional Parameters}) و دیگر متغیرهای ویژه سیستمی را نمی‌توان با استفاده از این روش (\cmd{:=}) مقداردهی کرد؛ چرا که این پارامترها ماهیتی «فقط‌خواندنی» (\en{Read-only}) دارند و مقدار آن‌ها مستقیماً توسط پوسته مدیریت می‌شود.
\end{note}

\textbf{\cmd{\$\{parameter:?word\}}}

چنانچه \file{parameter} تعریف‌نشده یا تهی باشد، این بسط منجر به توقف فوری اجرای اسکریپت شده و محتوای رشته \file{word} را به عنوان پیام خطا در خروجی خطای استاندارد (\en{stderr}) درج می‌کند. اما در صورتی که پارامتر دارای مقدار باشد، فرآیند بسط با موفقیت انجام شده و مقدار متغیر بازگردانده می‌شود.

\begin{latin}
\begin{lstlisting}
[me@linuxbox ~]$ foo=
[me@linuxbox ~]$ echo ${foo:?"parameter is empty"}
bash: foo: parameter is empty
[me@linuxbox ~]$ echo $?
1
[me@linuxbox ~]$ foo=bar
[me@linuxbox ~]$ echo ${foo:?"parameter is empty"}
bar
[me@linuxbox ~]$ echo $?
0
\end{lstlisting}
\end{latin}

این ساختار به‌ویژه برای اعتبارسنجی متغیرهای حیاتی در ابتدای اسکریپت‌ها بسیار مفید است؛ چرا که از ادامه کار برنامه با مقادیر نامعتبر جلوگیری می‌کند.

\textbf{\cmd{\$\{parameter:+word\}}}

این بسط دقیقاً عملکردی معکوس الگوهای قبلی دارد؛ به این صورت که اگر \file{parameter} تعریف‌نشده یا تهی باشد، هیچ خروجی تولید نمی‌کند. اما چنانچه پارامتر دارای مقدار باشد، رشته \file{word} جایگزین آن در خروجی می‌شود. نکته مهم این است که مقدار اصلی متغیر تحت هیچ شرایطی تغییر نمی‌کند.

\begin{latin}
\begin{lstlisting}
[me@linuxbox ~]$ foo=
[me@linuxbox ~]$ echo ${foo:+"substitute value if set"}

[me@linuxbox ~]$ foo=bar
[me@linuxbox ~]$ echo ${foo:+"substitute value if set"}
substitute value if set
\end{lstlisting}
\end{latin}

این ساختار عمدتاً برای بررسی وجود یک متغیر و اجرای عملیاتی مشروط به آن (مانند افزودن یک گزینه به دستور تنها در صورت وجود مقدار) بدون دست‌کاری محتوای اصلی متغیر کاربرد دارد.

\subsection{بسط‌های بازگرداننده نام متغیرها}

پوسته \en{Bash} از توانایی شناسایی و بازگرداندن نام متغیرهای تعریف‌شده برخوردار است. این قابلیت در سناریوهای برنامه‌نویسی خاص، به‌ویژه زمانی که نیاز به پیمایش در متغیرهای سیستم یا یافتن متغیرهایی با الگوی نام‌گذاری مشترک دارید، بسیار کاربردی است.

\textbf{\cmd{\$\{!prefix\*\}}}
\textbf{\cmd{\$\{!prefix@\}}}

این دو ساختار، فهرست نام تمامی متغیرهایی را که با عبارت \file{prefix} آغاز می‌شوند، بازمی‌گردانند. مطابق با مستندات رسمی \en{Bash}، هر دو فرمت عملکرد یکسانی دارند. در مثال زیر، تمامی متغیرهای محیطی که با رشته \en{BASH} شروع می‌شوند، فهرست شده‌اند:

\begin{latin}
\begin{lstlisting}
[me@linuxbox ~]$ echo ${!BASH*}
BASH BASH_ARGC BASH_ARGV BASH_COMMAND BASH_COMPLETION BASH_COMPLETION_DIR BASH_LINENO BASH_SOURCE BASH_SUBSHELL BASH_VERSINFO BASH_VERSION
\end{lstlisting}
\end{latin}

این ویژگی به‌ویژه زمانی مفید است که بخواهید مجموعه‌ای از متغیرهای مرتبط را که به صورت پویا ایجاد شده‌اند، در یک حلقه پردازش کنید.

\subsection{عملیات روی رشته‌ها}

مجموعه وسیعی از بسط‌های پارامتر برای پردازش رشته‌ها در \en{Bash} تعبیه شده است که بسیاری از آن‌ها به‌طور ویژه‌ای برای پردازش مسیر فایل‌ها (\en{Pathnames}) کارآمد هستند.

\textbf{\cmd{\$\{\#parameter\}}}

این بسط، طول رشته موجود در \file{parameter} را بازمی‌گرداند. در حالت معمول، خروجی تعداد کاراکترهای رشته است؛ اما اگر به جای نام متغیر از کاراکترهای \cmd{@} یا \cmd{*} استفاده شود، خروجی معادل تعداد پارامترهای موقعیتی (آرگومان‌های ورودی) خواهد بود.

\begin{latin}
\begin{lstlisting}
[me@linuxbox ~]$ foo="This string is long."
[me@linuxbox ~]$ echo "'$foo' is ${#foo} characters long."
'This string is long.' is 20 characters long.
\end{lstlisting}
\end{latin}

\textbf{\cmd{\$\{parameter:offset\}}}
\textbf{\cmd{\$\{parameter:offset:length\}}}

این بسط‌ها جهت استخراج زیررشته (\en{Substring Extraction}) استفاده می‌شوند. فرآیند استخراج از موقعیت آفست (\en{offset}) شروع شده و تا انتهای رشته ادامه می‌یابد، مگر آنکه طول مشخصی (\en{length}) تعیین شده باشد.

\begin{latin}
\begin{lstlisting}
[me@linuxbox ~]$ foo="This string is long."
[me@linuxbox ~]$ echo ${foo:5}
string is long.
[me@linuxbox ~]$ echo ${foo:5:6}
string
\end{lstlisting}
\end{latin}

چنانچه مقدار \file{offset} منفی باشد، عملیات استخراج از انتهای رشته محاسبه می‌شود. در این حالت، درج یک فاصله قبل از عدد منفی الزامی است تا پوسته آن را با ساختار \cmd{\$\{parameter:-word\}} (تخصیص مقدار پیش‌فرض) اشتباه نگیرد. همچنین توجه داشته باشید که مقدار \file{length} نباید کمتر از صفر باشد.

اگر از \cmd{@} به عنوان پارامتر استفاده شود، نتیجه بسط شامل تعداد \file{length} پارامتر موقعیتی خواهد بود که از شاخص \file{offset} آغاز می‌شوند.

\begin{latin}
\begin{lstlisting}
[me@linuxbox ~]$ foo="This string is long."
[me@linuxbox ~]$ echo ${foo: -5}
long.
[me@linuxbox ~]$ echo ${foo: -5:2}
lo
\end{lstlisting}
\end{latin}

\textbf{\cmd{\$\{parameter\#pattern\}}}
\textbf{\cmd{\$\{parameter\#\#pattern\}}}

این دو ساختار برای حذف پیشوند (\en{Prefix Removal}) بر اساس یک الگوی مشخص (\file{pattern}) استفاده می‌شوند. الگو در اینجا از نویسه‌های جانشین (\en{Wildcard}) مشابه با الگوهای بسط مسیر فایل (\en{Globbing}) پیروی می‌کند. تفاوت میان این دو در نحوه برخورد با تطبیق‌های چندگانه است: ساختار نخست (\cmd{\#}) کوتاه‌ترین تطبیق ممکن و ساختار دوم (\cmd{\#\#}) بلندترین تطبیق ممکن را از ابتدای رشته حذف می‌کند.

در مثال زیر، تأثیر این دو روش بر روی یک فایل با پسوند چندگانه مشاهده می‌شود:

\begin{latin}
\begin{lstlisting}
[me@linuxbox ~]$ foo=file.txt.zip
[me@linuxbox ~]$ echo ${foo#*.}
txt.zip
[me@linuxbox ~]$ echo ${foo##*.}
zip
\end{lstlisting}
\end{latin}

\textbf{\cmd{\$\{parameter\%pattern\}}}
\textbf{\cmd{\$\{parameter\%\%pattern\}}}

این دو ساختار دقیقاً مشابه الگوهای \cmd{\#} و \cmd{\#\#} عمل می‌کنند، با این تفاوت کلیدی که عملیات حذف را از انتهای رشته (پسوند) به سمت ابتدا انجام می‌دهند. در اینجا نیز الگوی اول (\cmd{\%}) کوتاه‌ترین تطبیق و الگوی دوم (\cmd{\%\%}) بلندترین تطبیق ممکن را از انتهای محتوای پارامتر حذف می‌کند.

مثال زیر نحوه جداسازی لایه‌های مختلف پسوند یک فایل را با این دو روش نشان می‌دهد:

\begin{latin}
\begin{lstlisting}
[me@linuxbox ~]$ foo=file.txt.zip
[me@linuxbox ~]$ echo ${foo%.*}
file.txt
[me@linuxbox ~]$ echo ${foo%%.*}
file
\end{lstlisting}
\end{latin}

\begin{itemize}
  \item \textbf{\cmd{\$\{parameter/pattern/string\}}}
  \item \textbf{\cmd{\$\{parameter//pattern/string\}}}
  \item \textbf{\cmd{\$\{parameter/\#pattern/string\}}}
  \item \textbf{\cmd{\$\{parameter/\%pattern/string\}}}
\end{itemize}

این ساختارها امکان عملیات جستجو و جایگزینی (\en{Search and Replace}) را مستقیماً بر روی محتوای یک پارامتر فراهم می‌کنند. چنانچه بخشی از رشته با الگوی تعیین‌شده (\file{pattern}) مطابقت داشته باشد، با عبارت \file{string} جایگزین خواهد شد.

تفاوت عملکرد در این چهار حالت به شرح زیر است:

\begin{itemize}
  \item \textbf{حالت اول (\cmd{/}):} تنها نخستین مورد انطباق یافته را جایگزین می‌کند.
  \item \textbf{حالت دوم (\cmd{//}):} تمام موارد انطباق یافته در کل رشته را جایگزین می‌کند.
  \item \textbf{حالت سوم (\cmd{/\#}):} جایگزینی را تنها در صورتی انجام می‌دهد که انطباق در ابتدای رشته رخ دهد.
  \item \textbf{حالت چهارم (\cmd{/\%}):} جایگزینی را تنها در صورتی انجام می‌دهد که انطباق در انتهای رشته رخ دهد.
\end{itemize}

\begin{note}
\textbf{نکته:} در تمام این حالات، در صورت حذف بخش \cmd{/string}، متن منطبق با الگو صرفاً حذف می‌شود.
\end{note}

مثال عملی برای درک تفاوت‌ها:

\begin{latin}
\begin{lstlisting}
[me@linuxbox ~]$ foo=JPG.JPG
[me@linuxbox ~]$ echo ${foo/JPG/jpg}
jpg.JPG
[me@linuxbox ~]$ echo ${foo//JPG/jpg}
jpg.jpg
[me@linuxbox ~]$ echo ${foo/#JPG/jpg}
jpg.JPG
[me@linuxbox ~]$ echo ${foo/%JPG/jpg}
JPG.jpg
\end{lstlisting}
\end{latin}

بسط پارامتر (\en{Parameter Expansion}) ابزاری قدرتمند در پوسته است که می‌تواند در بسیاری از موارد به عنوان جایگزینی بهینه برای دستورات خارجی نظیر \cmd{sed} یا \cmd{cut} عمل کند. استفاده از قابلیت‌های داخلی پوسته (\en{Built-in}) به جای فراخوانی برنامه‌های خارجی، با حذف سربار ناشی از ایجاد زیرپوسته (\en{Subshell})، کارایی و بهره‌وری اسکریپت‌ها را به شکل قابل توجهی افزایش می‌دهد.

به عنوان یک نمونه عملی، اسکریپت \cmd{longest-word} را که در فصل گذشته بررسی شد، بازنویسی می‌کنیم. در این نسخه، به جای استفاده از ساختار \cmd{\$(echo -n \$j | wc -c)}، از بسط پارامتر \cmd{\$\{\#j\}} برای محاسبه طول رشته استفاده شده است:

\begin{latin}
\begin{lstlisting}
#!/bin/bash

# longest-word3: find longest string in a file

for i; do
    if [[ -r "$i" ]]; then
        max_word=
        max_len=0
        for j in $(strings $i); do
            len="${#j}"
            if (( len > max_len )); then
                max_len="$len"
                max_word="$j"
            fi
        done
        echo "$i: '$max_word' ($max_len characters)"
    fi
done
\end{lstlisting}
\end{latin}

اکنون با استفاده از دستور \cmd{time}، تفاوت عملکرد بین نسخه اصلی (که از فراخوانی برنامه خارجی استفاده می‌کرد) و نسخه بهینه‌سازی شده را مقایسه می‌کنیم:

\begin{latin}
\begin{lstlisting}
[me@linuxbox ~]$ time longest-word2 dirlist-usr-bin.txt
dirlist-usr-bin.txt: 'scrollkeeper-get-extended-content-list' (38 characters)

real	0m3.618s
user	0m1.544s
sys	0m1.768s
[me@linuxbox ~]$ time longest-word3 dirlist-usr-bin.txt
dirlist-usr-bin.txt: 'scrollkeeper-get-extended-content-list' (38 characters)

real	0m0.060s
user	0m0.056s
sys	0m0.008s
\end{lstlisting}
\end{latin}

همان‌طور که مشاهده می‌شود، نسخه اصلی برای پردازش فایل متنی ۳.۶۱۸ ثانیه زمان صرف می‌کند، در حالی که نسخه جدید با بهره‌گیری از بسط پارامتر، تنها در ۰.۰۶ ثانیه عملیات را به پایان می‌رساند. این بهبود چشم‌گیر، اهمیت استفاده از قابلیت‌های داخلی پوسته را در اسکریپت‌نویسی حرفه‌ای لینوکس نشان می‌دهد.

\section{مکانیزم‌های تغییر حالت حروف (Case Conversion)}

پوسته \en{Bash} از چهار نوع بسط پارامتر و دو گزینه در دستور \cmd{declare} برای مدیریت و تبدیل حالت حروف (بزرگ/کوچک) در رشته‌ها پشتیبانی می‌کند.

اما اهمیت تغییر حالت حروف در چیست؟ فارغ از جنبه‌های زیبایی‌شناختی در نمایش خروجی، این قابلیت نقشی کلیدی در منطق برنامه‌نویسی ایفا می‌کند. سناریویی را در نظر بگیرید که در آن قصد جستجوی یک رشته در پایگاه داده را دارید. کاربر ممکن است داده را با حروف تماماً بزرگ، تماماً کوچک و یا ترکیبی از هر دو وارد کند. از آنجا که ذخیره تمام حالات ممکنِ ترکیب حروف در پایگاه داده غیرمنطقی است، از روش «نرمال‌سازی» (\en{Normalization}) استفاده می‌شود. در این روش، ورودی کاربر پیش از جستجو به یک فرم استاندارد (مثلاً تماماً کوچک) تبدیل می‌شود تا با داده‌های یکپارچه موجود در پایگاه داده مطابقت یابد.

استفاده از دستور \cmd{declare} یکی از روش‌های نرمال‌سازی رشته‌هاست. با این دستور می‌توان متغیری تعریف کرد که صرف‌نظر از مقدار ورودی، همواره فرمت مشخصی (تماماً بزرگ یا تماماً کوچک) را حفظ کند.

\begin{latin}
\begin{lstlisting}
#!/bin/bash

# ul-declare: demonstrate case conversion via declare

declare -u upper
declare -l lower

if [[ $1 ]]; then
    upper="$1"
    lower="$1"
    echo "$upper"
    echo "$lower"
fi
\end{lstlisting}
\end{latin}

در اسکریپت فوق، دو متغیر با نام‌های \cmd{upper} (با گزینه \cmd{-u} برای حروف بزرگ) و \cmd{lower} (با گزینه \cmd{-l} برای حروف کوچک) تعریف شده‌اند. پس از تخصیص اولین پارامتر موقعیتی به این متغیرها، خروجی به صورت خودکار نرمال‌سازی می‌شود:

\begin{latin}
\begin{lstlisting}
[me@linuxbox ~]$ ul-declare aBc
ABC
abc
\end{lstlisting}
\end{latin}

علاوه بر دستور \cmd{declare}، چهار الگوی بسط پارامتر نیز برای کنترل دقیق‌تر بر تبدیل حروف در دسترس است که در جدول زیر تشریح شده‌اند:

\begin{table}[htbp]
\centering
\caption{الگوهای بسط پارامتر برای تبدیل حالت حروف}
\label{tab:case-conversion}
\begin{tabularx}{\textwidth}{l X}
\toprule
\textbf{فرمت نگارش} & \textbf{شرح عملکرد} \\
\midrule
\cmd{\$\{parameter,,\}} یا \cmd{\$\{parameter,,pattern\}} & تمام حروف پارامتر را به حروف کوچک تبدیل می‌کند. عبارت \file{pattern} اختیاری است و برای محدود کردن دامنه‌ی تغییرات به نویسه‌های خاص (مانند \cmd{[A-F]}) به کار می‌رود. \\
\addlinespace
\cmd{\$\{parameter,pattern\}} & تنها نخستین نویسه‌ی پارامتر را به حروف کوچک تبدیل می‌کند. \\
\addlinespace
\cmd{\$\{parameter\^{}\^{}pattern\}} & تمام حروف پارامتر را به حروف بزرگ تبدیل می‌کند. \\
\addlinespace
\cmd{\$\{parameter\^{}pattern\}} & تنها نخستین نویسه‌ی پارامتر را به حروف بزرگ تبدیل می‌کند. \\
\bottomrule
\end{tabularx}
\end{table}

در ادامه، اسکریپتی جهت نمایش کاربردی این بسط‌ها ارائه شده است:

\begin{latin}
\begin{lstlisting}
#!/bin/bash

# ul-param: demonstrate case conversion via parameter expansion

if [[ "$1" ]]; then
    echo "${1,,}"
    echo "${1,}"
    echo "${1^^}"
    echo "${1^}"
fi
\end{lstlisting}
\end{latin}

نمونه‌ای از اجرای اسکریپت فوق را مشاهده می‌کنید:

\begin{latin}
\begin{lstlisting}
[me@linuxbox ~]$ ul-param aBc
abc
aBc
ABC
ABc
\end{lstlisting}
\end{latin}

در این مثال، اولین آرگومان ورودی (پارامتر موقعیتی یک) مورد پردازش قرار گرفته و هر چهار حالت عملیاتیِ بسط پارامتر بر روی آن اعمال شده است. شایان ذکر است که اگرچه در این نمونه از پارامتر موقعیتی استفاده شده، اما در عمل \file{parameter} می‌تواند هرگونه رشته ثابت، متغیر تعریف‌شده یا عبارت رشته‌ای معتبر باشد.

\section{ارزیابی و بسط محاسباتی}

در فصل ۷، به بررسی مفهوم بسط محاسباتی (\en{Arithmetic Expansion}) پرداختیم. این قابلیت جهت اجرای عملیات‌های ریاضی متنوع بر روی اعداد صحیح (\en{Integers}) طراحی شده است. ساختار پایه این بسط به صورت زیر است:

\begin{latin}
\begin{lstlisting}
$((expression))
\end{lstlisting}
\end{latin}

در این ساختار، \file{expression} بیانگر یک عبارت محاسباتی معتبر است.

این مفهوم با دستور مرکب \cmd{(( ))} که پیش‌تر در فصل ۲۷ برای ارزیابی محاسباتی (\en{Arithmetic Evaluation}) و شرط‌های منطقی با آن آشنا شدیم، قرابت نزدیکی دارد. در فصول قبلی، برخی از عملگرها و عبارات رایج را بررسی کردیم؛ اکنون به بررسی فهرست جامع‌تری از این قابلیت‌ها خواهیم پرداخت.

\subsection{مبانی سیستم‌های عددی در پوسته}

در فصل ۹، با مفاهیم مقدماتی اعداد در مبنای ۸ (هشت‌تایی) و مبنای ۱۶ (شانزده‌تایی) آشنا شدیم. در حوزه عبارات محاسباتی (\en{Arithmetic Expressions})، پوسته \en{Bash} انعطاف‌پذیری بالایی داشته و از اعداد صحیح در تمامی مبناهای عددی پشتیبانی می‌کند. جدول زیر قراردادهای نگارشی مورد استفاده برای تعیین مبنای اعداد را تشریح می‌کند.

\begin{table}[htbp]
\centering
\caption{نشانه‌گذاری مبنای عددی (\en{Radix}) در پوسته}
\label{tab:radix-notation}
\begin{tabularx}{\textwidth}{l X}
\toprule
\textbf{نشانه‌گذاری} & \textbf{شرح فنی} \\
\midrule
\cmd{number} & در حالت پیش‌فرض، اعدادی که فاقد پیشوند هستند، در مبنای ۱۰ (ده‌دهی) ارزیابی می‌شوند. \\
\addlinespace
\cmd{0number} & در عبارات محاسباتی، اعدادی که با رقم صفر آغاز می‌شوند، به‌عنوان عدد در مبنای ۸ (\en{Octal}) تلقی می‌گردند. \\
\addlinespace
\cmd{0xnumber} & این پیشوند نشان‌دهنده‌ی عدد در مبنای ۱۶ (\en{Hexadecimal}) است. \\
\addlinespace
\cmd{base\#number} & در این ساختار، عدد \file{number} در مبنای \file{base} تفسیر می‌شود. \\
\bottomrule
\end{tabularx}
\end{table}

استفاده از ساختار \cmd{base\#number} به شما این امکان را می‌دهد که از هر مبنایی بین ۲ تا ۶۴ استفاده کنید. این قابلیت به‌ویژه در پردازش داده‌های سطح پایین، تنظیمات شبکه یا کار با کدهای کنترلی سخت‌افزار بسیار حائز اهمیت است.

در ادامه، نمونه‌هایی از نحوه به‌کارگیری این نشانه‌گذاری‌ها ارائه شده است:

\begin{latin}
\begin{lstlisting}
[me@linuxbox ~]$ echo $((0xff))
255
[me@linuxbox ~]$ echo $((2#11111111))
255
\end{lstlisting}
\end{latin}

در مثال‌های فوق، ابتدا مقدار معادل عدد هگزادسیمال \cmd{ff} (بزرگ‌ترین مقدار قابل نمایش در دو رقمِ مبنای ۱۶) و سپس معادل بزرگ‌ترین عدد دودویی ۸ رقمی (مبنای ۲) در سیستم ده‌دهی چاپ شده است. این خروجی‌ها نشان می‌دهند که چگونه پوسته به‌صورت خودکار تمامی مبناهای ورودی را جهت نمایش یا محاسبات بعدی به مبنای ۱۰ تبدیل می‌کند.

\subsection{عملگرهای یگانی (Unary Operators)}

در محاسبات پوسته، دو عملگر یگانی \cmd{+} و \cmd{-} جهت تعیین علامت عدد به کار می‌روند. این عملگرها مستقیماً پیش از مقدار عددی قرار می‌گیرند؛ به عنوان نمونه، \cmd{-5} نشان‌دهنده یک مقدار صحیح منفی است.

\subsection{محاسبات پایه (Simple Arithmetic)}

عملگرهای استاندارد ریاضی که در پردازش‌های محاسباتی \en{Bash} مورد استفاده قرار می‌گیرند، در جدول زیر درج شده‌اند. توجه داشته باشید که این عملگرها در ساختار بسط محاسباتی، صرفاً بر روی اعداد صحیح عمل می‌کنند.

\begin{table}[htbp]
\centering
\caption{عملگرهای محاسباتی}
\label{tab:arith-ops-basic}
\begin{tabularx}{\textwidth}{c X}
\toprule
\textbf{عملگر} & \textbf{شرح عملکرد} \\
\midrule
\cmd{+} & جمع \\
\addlinespace
\cmd{-} & تفریق \\
\addlinespace
\cmd{*} & ضرب \\
\addlinespace
\cmd{/} & تقسیم صحیح (بخش اعشاری نتیجه حذف می‌شود) \\
\addlinespace
\cmd{**} & توان \\
\addlinespace
\cmd{\%} & باقیمانده (باقی‌ماندهٔ حاصل از تقسیم دو عدد بر یکدیگر) \\
\bottomrule
\end{tabularx}
\end{table}

در محاسبات پوسته، عملگر تقسیم (\cmd{/}) به صورت «تقسیم صحیح» عمل می‌کند؛ به این معنا که اگر نتیجه تقسیم دارای بخش اعشاری باشد، آن بخش نادیده گرفته شده و تنها جزء صحیح بازگردانده می‌شود. همچنین عملگر \cmd{\%} (باقی‌مانده (\en{Modulo})) برای به‌دست آوردن باقیمانده تقسیم بسیار کاربردی است.

بسیاری از این عملگرها ماهیتی بدیهی دارند، اما مفاهیم تقسیم صحیح (\en{Integer Division}) و باقیمانده (\en{Modulo}) نیازمند تبیین دقیق‌تری هستند.

\begin{latin}
\begin{lstlisting}
[me@linuxbox ~]$ echo $(( 5 / 2 ))
2
\end{lstlisting}
\end{latin}

در چنین سیستمی، تعیین باقیمانده تقسیم اهمیت دوچندانی پیدا می‌کند؛ چرا که مکمل عملیات تقسیم صحیح است:

\begin{latin}
\begin{lstlisting}
[me@linuxbox ~]$ echo $(( 5 % 2 ))
1
\end{lstlisting}
\end{latin}

با تلفیق این دو عملگر، می‌توان دریافت که در تقسیم عدد ۵ بر ۲، خارج‌قسمت برابر با ۲ و باقیمانده معادل ۱ است.

محاسبه باقیمانده در ساختارهای تکرار (حلقه‌ها) بسیار کلیدی است؛ این قابلیت به برنامه‌نویس اجازه می‌دهد تا عملیات خاصی را در فواصل زمانی یا گام‌های مشخص اجرا کند. در اسکریپت زیر، فهرستی از اعداد چاپ شده و با استفاده از عملگر باقی‌مانده (\en{Modulo})، مضارب عدد ۵ متمایز می‌شوند:

\begin{latin}
\begin{lstlisting}
#!/bin/bash

# modulo: demonstrate the modulo operator

for ((i = 0; i <= 20; i = i + 1)); do
    remainder=$((i % 5))
    if (( remainder == 0 )); then
        printf "<%d> " "$i"
    else
        printf "%d " "$i"
    fi
done
printf "\n"
\end{lstlisting}
\end{latin}

خروجی حاصل از اجرای این اسکریپت به شرح زیر است:

\begin{latin}
\begin{lstlisting}
[me@linuxbox ~]$ modulo
<0> 1 2 3 4 <5> 6 7 8 9 <10> 11 12 13 14 <15> 16 17 18 19 <20>
\end{lstlisting}
\end{latin}

\subsection{عملگرهای انتساب (Assignment)}

اگرچه ممکن است کاربرد مستقیم انتساب در عبارات محاسباتی (\en{Arithmetic Expression}) در ابتدا چندان ملموس به نظر نرسد، اما پوسته \en{Bash} اجازه می‌دهد تا این فرآیند را درون ساختارهای محاسباتی به انجام برسانیم. ما پیش‌تر به‌صورت معمول مقادیر را به متغیرها تخصیص می‌دادیم، اما اجرای این عمل درون یک عبارت محاسباتی، قابلیت‌های جدیدی را فراهم می‌کند.

\begin{latin}
\begin{lstlisting}
[me@linuxbox ~]$ foo=
[me@linuxbox ~]$ echo $foo

[me@linuxbox ~]$ if (( foo = 5 )); then echo "It is true."; fi
It is true.
[me@linuxbox ~]$ echo $foo
5
\end{lstlisting}
\end{latin}

در مثال فوق، ابتدا متغیر \cmd{foo} را تهی کرده و وضعیت آن را بررسی می‌کنیم. سپس با استفاده از دستور مرکب \cmd{(( foo = 5 ))} در ساختار شرطی \cmd{if}، دو رخداد هم‌زمان شکل می‌گیرد: نخست، مقدار ۵ به متغیر \cmd{foo} تخصیص می‌یابد و دوم، از آنجا که خروجی این انتساب عددی غیرصفر است، عبارت از نظر منطقی «درست» (\en{true}) ارزیابی می‌شود.

\begin{note}
\textbf{نکته:} درک تمایز دقیق میان عملگر \cmd{=} و \cmd{==} در این عبارات حیاتی است. عملگر \cmd{=} منحصراً برای «انتساب» به کار می‌رود؛ عبارت \cmd{foo = 5} به معنای «مقدار ۵ را در متغیر foo جای‌گذاری کن» است. در مقابل، عملگر \cmd{==} برای «سنجش برابری» استفاده می‌شود؛ عبارت \cmd{foo == 5} یعنی «آیا مقدار فعلی متغیر foo برابر با ۵ است؟». این تفاوت در بسیاری از زبان‌های برنامه‌نویسی وجود دارد، اما در پوسته \en{Bash} به دلیل استفاده از تک‌علامت \cmd{=} برای مقایسه رشته‌ها در دستور \cmd{test}، ممکن است موجب سردرگمی شود. به همین دلیل، استفاده از دستورات مدرن‌تر نظیر \cmd{[[ ]]} و \cmd{(( ))} به جای \cmd{test} توصیه می‌شود.
\end{note}

علاوه بر انتساب ساده، پوسته نشانه‌گذاری‌های ترکیبی و کاربردی دیگری را نیز ارائه می‌دهد که در جدول زیر فهرست شده‌اند:

\begin{table}[htbp]
\centering
\caption{عملگرهای انتساب و تغییر مقدار}
\label{tab:assignment-ops}
\begin{tabularx}{\textwidth}{l X}
\toprule
\textbf{نشانه‌گذاری} & \textbf{توضیح} \\
\midrule
\cmd{parameter = value} & انتساب ساده: مقدار \file{value} به پارامتر اختصاص می‌یابد. \\
\addlinespace
\cmd{parameter += value} & انتساب جمع: معادل \cmd{parameter = parameter + value} \\
\addlinespace
\cmd{parameter -= value} & انتساب تفریق: معادل \cmd{parameter = parameter - value} \\
\addlinespace
\cmd{parameter *= value} & انتساب ضرب: معادل \cmd{parameter = parameter * value} \\
\addlinespace
\cmd{parameter /= value} & انتساب تقسیم صحیح: معادل \cmd{parameter = parameter / value} \\
\addlinespace
\cmd{parameter \%= value} & انتساب باقیمانده: معادل \cmd{parameter = parameter \% value} \\
\addlinespace
\cmd{parameter++} & پس‌افزایش (\en{Post-increment}): مقدار متغیر را پس از ارزیابی عبارت، یک واحد افزایش می‌دهد. \\
\addlinespace
\cmd{parameter--} & پس‌کاهش (\en{Post-decrement}): مقدار متغیر را پس از ارزیابی عبارت، یک واحد کاهش می‌دهد. \\
\addlinespace
\cmd{++parameter} & پیش‌افزایش (\en{Pre-increment}): مقدار متغیر را پیش از ارزیابی عبارت، یک واحد افزایش می‌دهد. \\
\addlinespace
\cmd{--parameter} & پیش‌کاهش (\en{Pre-decrement}): مقدار متغیر را پیش از ارزیابی عبارت، یک واحد کاهش می‌دهد. \\
\bottomrule
\end{tabularx}
\end{table}

این عملگرهای انتساب، ساختاری فشرده و کارآمد برای اجرای بسیاری از وظایف متداول محاسباتی فراهم می‌آورند. در این میان، عملگرهای افزایش (\cmd{++}) و کاهش (\cmd{--}) به دلیل توانایی در تغییر مقدار پارامتر به میزان یک واحد، اهمیت ویژه‌ای دارند. این شیوه نگارش که از زبان برنامه‌نویسی \en{C} اقتباس شده، در بسیاری از زبان‌های دیگر و از جمله \en{Bash} به کار گرفته شده است.

این عملگرها می‌توانند به صورت پیشوند (قبل از نام متغیر) یا پسوند (بعد از نام متغیر) ظاهر شوند. با وجود اینکه هر دو حالت منجر به تغییر مقدار متغیر می‌شوند، اما در زمان بازگشت مقدار، تفاوت ظریفی میان آن‌ها وجود دارد: در حالت پیشوند، ابتدا تغییر مقدار اعمال شده و سپس مقدار جدید بازگردانده می‌شود. در حالت پسوند، ابتدا مقدار فعلی بازگردانده شده و سپس عمل افزایش یا کاهش صورت می‌گیرد. این رفتار اگرچه در ابتدا نامتعارف به نظر می‌رسد، اما در منطق برنامه‌نویسی کاربردهای خاص خود را دارد.

نمونه زیر این تفاوت عملکردی را به خوبی نشان می‌دهد:

\begin{latin}
\begin{lstlisting}
[me@linuxbox ~]$ foo=1
[me@linuxbox ~]$ echo $((foo++))
1
[me@linuxbox ~]$ echo $foo
2
\end{lstlisting}
\end{latin}

در مثال بالا، با استفاده از عملگر پسوند، مقدار اولیه (\cmd{1}) چاپ می‌شود، اما در فراخوانی بعدی مشاهده می‌کنیم که مقدار به \cmd{2} افزایش یافته است. در مقابل، استفاده از عملگر پیشوند منجر به رفتاری مستقیم‌تر می‌شود:

\begin{latin}
\begin{lstlisting}
[me@linuxbox ~]$ foo=1
[me@linuxbox ~]$ echo $((++foo))
2
[me@linuxbox ~]$ echo $foo
2
\end{lstlisting}
\end{latin}

در اکثر سناریوهای اسکریپت‌نویسی پوسته، استفاده از حالت پیشوند معمولاً به نتایج قابل‌انتظارتری منجر می‌شود. عملگرهای \cmd{++} و \cmd{--} غالباً در ساختار حلقه‌ها به کار می‌روند. در ادامه، نسخه بهینه‌شده‌ای از اسکریپت قبلی (محاسبه باقیمانده) ارائه شده است که از نوشتار فشرده‌تری بهره می‌برد:

\begin{latin}
\begin{lstlisting}
#!/bin/bash

# modulo2: demonstrate the modulo operator

for ((i = 0; i <= 20; ++i)); do
    if (((i % 5) == 0)); then
        printf "<%d> " "$i"
    else
        printf "%d " "$i"
    fi
done
printf "\n"
\end{lstlisting}
\end{latin}

\subsection{عملیات بیتی (Bitwise Operations)}

دسته‌ای از عملگرها در پوسته \en{Bash} وجود دارند که به شیوه‌ای متمایز با اعداد برخورد می‌کنند. این عملگرها مستقیماً بر روی سطح بیت‌ها عمل کرده و غالباً برای وظایف سطح پایین (\en{Low-level})، نظیر تنظیم (\en{Set}) یا خواندن پرچم‌های بیتی (\en{Bit-flags}) و ماسک‌های شبکه (\en{Network Masks}) استفاده می‌شوند. عملگرهای بیتی در جدول زیر تشریح شده‌اند:

\begin{table}[htbp]
\centering
\caption{عملگرهای بیتی}
\label{tab:bitwise-ops}
\begin{tabularx}{\textwidth}{c X}
\toprule
\textbf{عملگر} & \textbf{شرح عملکرد فنی} \\
\midrule
\cmd{\textasciitilde} & نقیض بیتی (\en{Bitwise NOT}): تمامی بیت‌های عدد را معکوس می‌کند (0 به 1 و 1 به 0). \\
\addlinespace
\cmd{<<} & شیفت به چپ (\en{Left Shift}): تمامی بیت‌های عدد اول را به تعداد مشخص‌شده توسط عدد دوم، به سمت چپ جابه‌جا می‌کند و جایگاه‌های خالی سمت راست را با صفر پر می‌کند؛ این عمل معادل ضرب عدد در توان‌های ۲ است. \\
\addlinespace
\cmd{>>} & شیفت به راست (\en{Right Shift}): تمامی بیت‌های عدد اول را به تعداد مشخص‌شده توسط عدد دوم، به سمت راست جابه‌جا می‌کند؛ این عمل معادل تقسیم صحیح عدد بر توان‌های ۲ است. \\
\addlinespace
\cmd{\&} & \en{AND} بیتی (\en{Bitwise AND}): عملگر دوتایی است که عملیات منطقی \en{AND} را بر روی هر جفت بیت متناظر از دو عدد اعمال می‌کند؛ حاصل هر بیت تنها در صورتی 1 خواهد بود که هر دو بیت متناظر 1 باشند. \\
\addlinespace
\cmd{|} & \en{OR} بیتی (\en{Bitwise OR}): عملگر دوتایی است که عملیات منطقی \en{OR} را بر روی هر جفت بیت متناظر از دو عدد اعمال می‌کند؛ حاصل هر بیت در صورتی 1 خواهد بود که حداقل یکی از دو بیت متناظر ۱ باشد. \\
\addlinespace
\cmd{\^{}} & \en{XOR} بیتی (\en{Bitwise XOR / OR} انحصاری): عملگر دوتایی است که عملیات \en{OR} انحصاری را بر روی هر جفت بیت متناظر اعمال می‌کند؛ حاصل هر بیت تنها در صورتی 1 خواهد بود که دقیقاً یکی از دو بیت متناظر (و نه هر دو) برابر 1 باشد. \\
\bottomrule
\end{tabularx}
\end{table}

شایان ذکر است که برای تمامی این موارد (به‌جز نقیض بیتی)، عملگرهای انتساب متناظر (مانند \cmd{>>=}) نیز در دسترس هستند.

یکی از کاربردهای کلاسیک شیفت بیتی، انجام سریع محاسبات توان ۲ است. در مثال زیر، با استفاده از عملگر شیفت به چپ، فهرستی از توان‌های عدد ۲ تولید شده است:

\begin{latin}
\begin{lstlisting}
[me@linuxbox ~]$ for ((i=0;i<8;++i)); do echo $((1<<i)); done
1
2
4
8
16
32
64
128
\end{lstlisting}
\end{latin}

در این مثال، هر بار که بیت \cmd{1} یک واحد به سمت چپ شیفت داده می‌شود، مقدار ده‌دهی آن دو برابر می‌گردد که روشی بسیار بهینه برای تولید توالی توان‌های ۲ در اسکریپت‌نویسی است.

\subsection{عملگرهای مقایسه‌ای و منطقی}

همان‌طور که در فصل ۲۷ بررسی شد، دستور مرکب \cmd{(( ))} از طیف گسترده‌ای از عملگرهای مقایسه پشتیبانی می‌کند. علاوه بر این، چندین عملگر جهت ارزیابی‌های منطقی پیچیده‌تر نیز در دسترس هستند که در جدول زیر فهرست شده‌اند.

\begin{table}[htbp]
\centering
\caption{عملگرهای مقایسه‌ای و منطقی}
\label{tab:comparison-logical-ops}
\begin{tabularx}{\textwidth}{l X}
\toprule
\textbf{عملگر} & \textbf{شرح عملکرد فنی} \\
\midrule
\cmd{<=} & کوچک‌تر یا مساوی \\
\addlinespace
\cmd{>=} & بزرگ‌تر یا مساوی \\
\addlinespace
\cmd{<} & کوچک‌تر \\
\addlinespace
\cmd{>} & بزرگ‌تر \\
\addlinespace
\cmd{==} & برابری (\en{Equality}) \\
\addlinespace
\cmd{!=} & عدم برابری (\en{Inequality}) \\
\addlinespace
\cmd{\&\&} & \en{AND} منطقی (عملگر دوتایی): ترکیب دو عبارت شرطی، به‌گونه‌ای که حاصل تنها در صورتی صحیح (\en{True}) است که هر دو عبارت صحیح باشند. \\
\addlinespace
\cmd{||} & \en{OR} منطقی (عملگر دوتایی): ترکیب دو عبارت شرطی، به‌گونه‌ای که حاصل در صورتی صحیح (\en{True}) است که حداقل یکی از دو عبارت صحیح باشد. \\
\addlinespace
\cmd{expr1?expr2:expr3} & عملگر شرطی سه‌گانه (\en{Ternary Operator}): در صورتی که \file{expr1} صحیح (غیرصفر) ارزیابی شود، حاصل کل عبارت برابر با \file{expr2} خواهد بود؛ در غیر این صورت، حاصل برابر با \file{expr3} خواهد بود. \\
\bottomrule
\end{tabularx}
\end{table}

در محاسبات منطقی پوسته، عبارات از قواعد منطق ریاضی پیروی می‌کنند؛ به این معنا که مقدار صفر به عنوان «نادرست» (\en{false}) و هر مقدار غیرصفر به عنوان «درست» (\en{true}) تلقی می‌شود. دستور مرکب \cmd{(( ))} این نتایج محاسباتی را به کدهای وضعیت خروجی (\en{Exit Status}) استاندارد پوسته نگاشت می‌کند:

\begin{latin}
\begin{lstlisting}
[me@linuxbox ~]$ if ((1)); then echo "true"; else echo "false"; fi
true
[me@linuxbox ~]$ if ((0)); then echo "true"; else echo "false"; fi
false
\end{lstlisting}
\end{latin}

در میان این عملگرها، عملگر شرطی سه‌گانه (\en{Ternary Operator}) متمایزترین ساختار را دارد. این عملگر که از زبان برنامه‌نویسی \en{C} اقتباس شده، مانند یک نسخه فشرده از ساختار \cmd{if/then/else} عمل می‌کند. این عملگر بر روی سه عبارت محاسباتی اعمال می‌شود (توجه داشته باشید که بر روی رشته‌ها عمل نمی‌کند)؛ چنانچه عبارت اول (\file{expr1}) صحیح ارزیابی شود، عبارت دوم اجرا می‌گردد، در غیر این صورت عبارت سوم به اجرا در می‌آید.

مثال زیر نحوه عملکرد این عملگر را در خط فرمان نشان می‌دهد:

\begin{latin}
\begin{lstlisting}
[me@linuxbox ~]$ a=0
[me@linuxbox ~]$ ((a<1?++a:--a))
[me@linuxbox ~]$ echo $a
1
[me@linuxbox ~]$ ((a<1?++a:--a))
[me@linuxbox ~]$ echo $a
0
\end{lstlisting}
\end{latin}

در این نمونه، عملگر شرطی نقش یک سوییچ (\en{Toggle}) را ایفا می‌کند. با هر بار اجرای دستور، وضعیت متغیر \cmd{a} بین صفر و یک تغییر می‌کند که تکنیکی کارآمد برای مدیریت وضعیت‌های دوگانه در اسکریپت‌نویسی است.

باید توجه داشت که انجام عملیات انتساب به صورت مستقیم عملگر شرطی سه‌گانه (\en{Ternary Operator}) می‌تواند منجر به خطای نحوی شود:

\begin{latin}
\begin{lstlisting}
[me@linuxbox ~]$ a=0
[me@linuxbox ~]$ ((a<1?a+=1:a-=1))
bash: ((: a<1?a+=1:a-=1: attempted assignment to non-variable (error token is "-=1")
\end{lstlisting}
\end{latin}

این محدودیت به دلیل اولویت عملگرها ایجاد می‌شود و برای رفع آن، باید عبارات انتساب را درون پرانتز محصور کرد تا ترتیب اجرا به درستی رعایت شود:

\begin{latin}
\begin{lstlisting}
[me@linuxbox ~]$ ((a<1?(a+=1):(a-=1)))
\end{lstlisting}
\end{latin}

در ادامه، یک نمونه کامل از به‌کارگیری عملگرهای محاسباتی در قالب اسکریپتی که جدولی از توان‌های اعداد تولید می‌کند، آورده شده است:

\begin{latin}
\begin{lstlisting}
#!/bin/bash

# arith-loop: script to demonstrate arithmetic operators

finished=0
a=0
printf "a\ta**2\ta**3\n"
printf "=\t====\t====\n"

until ((finished)); do
    b=$((a**2))
    c=$((a**3))
    printf "%d\t%d\t%d\n" "$a" "$b" "$c"
    ((a<10?++a:(finished=1)))
done
\end{lstlisting}
\end{latin}

در این اسکریپت، از حلقه \cmd{until} بر مبنای وضعیت خروج (\en{Exit Status}) متغیر \cmd{finished} استفاده شده است. در ابتدا، مقدار این متغیر صفر است (که در منطق محاسباتی معادل \en{false} تلقی می‌شود)؛ بنابراین حلقه تا زمانی که این مقدار به غیرصفر تغییر نکند، به اجرای خود ادامه می‌دهد.

درون بدنه حلقه، مربع و مکعبِ متغیر شمارنده \cmd{a} محاسبه می‌گردد. در گام نهایی، مقدار شمارنده توسط عملگر شرطی ارزیابی می‌شود: اگر \cmd{a} کوچک‌تر از \cmd{10} باشد، یک واحد به آن افزوده می‌شود (پیش‌افزایش)؛ در غیر این صورت، مقدار متغیر \cmd{finished} برابر با \cmd{1} قرار می‌گیرد تا شرط خروج از حلقه (تبدیل به \en{true} حسابی) محقق شده و فرآیند خاتمه یابد.

خروجی حاصل از اجرای اسکریپت فوق به شرح زیر است:

\begin{latin}
\begin{lstlisting}
[me@linuxbox ~]$ arith-loop
a	a**2	a**3
=	====	====
0	0	0
1	1	1
2	4	8
3	9	27
4	16	64
5	25	125
6	36	216
7	49	343
8	64	512
9	81	729
10	100	1000
\end{lstlisting}
\end{latin}

این خروجی نشان‌دهنده دقت و سرعت محاسبات ریاضی در پوسته \en{Bash} است. با وجود قدرت بالای این ابزار در مدیریت اعداد صحیح، محدودیت اصلی آن عدم پشتیبانی بومی از محاسبات اعشاری است.

\section{ابزار bc: زبان محاسباتی با دقت دلخواه}

پیش‌تر مشاهده کردیم که پوسته \en{Bash} به خوبی از پس محاسبات بر پایه اعداد صحیح برمی‌آید؛ اما در مواجهه با محاسبات پیشرفته ریاضی یا کار با اعداد ممیز شناور (\en{Floating-point})، توانمندی مستقیم پوسته به پایان می‌رسد. در واقع، پوسته به‌صورت بومی از اعداد اعشاری پشتیبانی نمی‌کند. برای عبور از این محدودیت، ناگزیر به استفاده از ابزارهای خارجی هستیم. اگرچه بهره‌گیری از قابلیت‌های داخلی زبان‌هایی نظیر \en{Perl} یا \en{AWK} راهکار ممکنی است، اما بررسی آن‌ها در چارچوب این کتاب نمی‌گنجد.

رویکرد استاندارد دیگر، استفاده از ماشین‌حساب‌های تخصصی خط فرمان است. یکی از قدرتمندترین این ابزارها که در اکثر توزیع‌های لینوکس به‌صورت پیش‌فرض موجود است، برنامه \cmd{bc} نام دارد.

ابزار \cmd{bc} در واقع مفسر زبانی با نحو (\en{Syntax}) مشابه زبان \en{C} است که فایل‌های ورودی یا داده‌های ارسال شده به ورودی استاندارد (\en{stdin}) را پردازش و اجرا می‌کند. زبان محاسباتی \cmd{bc} از قابلیت‌های متعددی همچون تعریف متغیرها، ساختارهای کنترلی (حلقه‌ها) و توابع اختصاصی پشتیبانی می‌کند. در اینجا قصد نداریم تمام جزئیات این ابزار را واکاوی کنیم و صرفاً به معرفی اصول اولیه آن می‌پردازیم؛ چرا که مستندات کامل آن در صفحات راهنما (\cmd{man bc}) به‌خوبی تدوین شده است.

برای شروع، یک نمونه بسیار ساده را در نظر بگیرید. فرض کنید می‌خواهیم اسکریپتی برای محاسبه حاصل‌جمع \cmd{2 + 2} بنویسیم:

\begin{latin}
\begin{lstlisting}
/* A very simple bc script */

2 + 2
\end{lstlisting}
\end{latin}

خط نخست این کد، یک «توضیح» (\en{Comment}) است. ابزار \cmd{bc} در نگارش توضیحات کاملاً از نحو زبان \en{C} پیروی می‌کند؛ به این معنا که توضیحات (که می‌توانند چندین خط را در بر بگیرند) با کاراکترهای \cmd{/*} آغاز شده و با \cmd{*/} پایان می‌یابند.

\subsection{تعامل با bc}

چنانچه اسکریپت قبلی را در فایلی با نام \file{foo.bc} ذخیره کنیم، فرآیند اجرای آن به صورت زیر خواهد بود:

\begin{latin}
\begin{lstlisting}
[me@linuxbox ~]$ bc foo.bc
bc 1.06.94
Copyright 1991-1994, 1997, 1998, 2000, 2004, 2006 Free Software Foundation, Inc.
This is free software with ABSOLUTELY NO WARRANTY.
For details type `warranty'.
4
\end{lstlisting}
\end{latin}

با دقت در خروجی، نتیجه محاسبه را در خط پایانی و پس از بیانیه حق نشر (\en{Copyright}) مشاهده می‌کنیم. برای حذف این پیام‌های آغازین و دریافت خروجی خالص، می‌توان از گزینه \cmd{-q} (مخفف \en{quiet}) استفاده کرد.

ابزار \cmd{bc} علاوه بر اجرای فایل، قابلیت استفاده به صورت تعاملی (\en{Interactive}) را نیز دارد:

\begin{latin}
\begin{lstlisting}
[me@linuxbox ~]$ bc -q
2 + 2
4
quit
\end{lstlisting}
\end{latin}

در حالت تعاملی، با وارد کردن هر عبارت محاسباتی، نتیجه بلافاصله در خط بعدی نمایش داده می‌شود. درج دستور \cmd{quit} به جلسه تعاملی پایان می‌دهد.

همچنین، می‌توان اسکریپت‌ها یا عبارات را از طریق ورودی استاندارد (\en{stdin}) به \cmd{bc} هدایت کرد:

\begin{latin}
\begin{lstlisting}
[me@linuxbox ~]$ bc < foo.bc
4
\end{lstlisting}
\end{latin}

بهره‌گیری از ورودی استاندارد به این معناست که می‌توانیم از مکانیزم‌هایی نظیر \en{Here Documents}، \en{Here Strings} و پایپ‌لاین‌ها (\en{Pipes}) برای ارسال مستقیم دستورات استفاده کنیم. در نمونه زیر، از یک \en{here string} استفاده شده است:

\begin{latin}
\begin{lstlisting}
[me@linuxbox ~]$ bc <<< "2+2"
4
\end{lstlisting}
\end{latin}

\subsection{سناریوی عملی: محاسبه اقساط وام}

به عنوان یک مثال کاربردی در دنیای واقعی، اسکریپتی طراحی می‌کنیم که یکی از محاسبات متداول مالی، یعنی تعیین مبلغ اقساط ماهانه وام را انجام می‌دهد. در این اسکریپت، از مکانیزم \en{Here Document} برای ارسال مستقیم دستورات محاسباتی به ابزار \cmd{bc} بهره گرفته شده است:

\begin{latin}
\begin{lstlisting}
#!/bin/bash

# loan-calc: script to calculate monthly loan payments

PROGNAME="${0##*/}" # Use parameter expansion to get basename

usage () {
    cat << _EOF_
Usage: $PROGNAME PRINCIPAL INTEREST MONTHS

Where:
PRINCIPAL is the amount of the loan.
INTEREST is the APR as a number (7% = 0.07).
MONTHS is the length of the loan's term.

_EOF_
}

if (($# != 3)); then
    usage
    exit 1
fi

principal=$1
interest=$2
months=$3

bc <<- _EOF_
    scale = 10
    i = $interest / 12
    p = $principal
    n = $months
    a = p * ((i * ((1 + i) ^ n)) / (((1 + i) ^ n) - 1))
    print a, "\n"
_EOF_
\end{lstlisting}
\end{latin}

پس از اجرا، خروجی به صورت زیر خواهد بود:

\begin{latin}
\begin{lstlisting}
[me@linuxbox ~]$ loan-calc 135000 0.0775 180
1270.7222490000
\end{lstlisting}
\end{latin}

در این مثال، قسط ماهانه یک وام ۱۳۵,۰۰۰ دلاری با نرخ بهره سالانه ۷.۷۵٪ در یک دوره ۱۸۰ ماهه (۱۵ سال) محاسبه شده است. دقت بالای خروجی مدیون متغیر ویژه \cmd{scale} در محیط \cmd{bc} است که تعداد ارقام اعشار را تعیین می‌کند.

زبان اسکریپت‌نویسی \cmd{bc} در صفحات راهنما (\cmd{man bc}) به‌طور کامل مستند شده است. اگرچه قواعد نگارش ریاضی آن تفاوت‌های اندکی با پوسته \en{Bash} دارد (و بیشتر به ساختار زبان \en{C} متمایل است)، اما بر اساس آموخته‌های پیشین، اکثر بخش‌های آن برای شما ملموس و آشنا خواهد بود.

\section{جمع‌بندی}

در این فصل، با مجموعه‌ای از جزئیات فنی و کاربردی آشنا شدیم که سنگ‌بنای انجام «عملیات واقعی» در اسکریپت‌نویسی محسوب می‌شوند. با ارتقای سطح مهارت شما در نگارش اسکریپت، توانایی پردازش بهینه رشته‌ها و مدیریت دقیق محاسبات عددی، ابزاری بسیار ارزشمند خواهد بود. اسکریپت \cmd{loan-calc} به‌خوبی نشان داد که چگونه حتی با قطعه کدهایی ساده، می‌توان ابزارهایی کاربردی و گره‌گشا خلق کرد.

\section{تمرین‌های تکمیلی}

اگرچه عملکرد هسته مرکزی اسکریپت \cmd{loan-calc} پیاده‌سازی شده است، اما این برنامه همچنان فضای زیادی برای بهبود دارد. برای ارتقای سطح تخصصی خود، سعی کنید ویژگی‌های زیر را به اسکریپت اضافه کنید:

\begin{itemize}
  \item \textbf{اعتبارسنجی ورودی‌ها:} پیاده‌سازی مکانیزمی جهت بررسی دقیق و کامل آرگومان‌های ورودی خط فرمان (مانند اطمینان از عددی بودن مقادیر).
  \item \textbf{حالت تعاملی (\en{Interactive Mode}):} افزودن یک گزینه به خط فرمان که در صورت فعال‌سازی، مقادیر وام، نرخ بهره و مدت‌زمان را به‌صورت پرسش‌وپاسخ از کاربر دریافت کند.
  \item \textbf{بهبود بصری خروجی:} استفاده از ابزارهایی نظیر \cmd{printf} برای قالب‌بندی شکیل‌تر و خواناتر نتایج نهایی.
\end{itemize}

\section{منابع مطالعه تکمیلی}

راهنمای مرجع \en{Bash} بحث جامعی درباره بسط پارامترها (\en{parameter expansion}) ارائه می‌دهد:

\begin{latin}
\url{http://www.gnu.org/software/bash/manual/bashref.html#Shell-Parameter-Expansion}
\end{latin}

ویکی‌پدیا مقاله‌ای مفید در توصیف عملیات بیتی (\en{bit operations}) دارد:

\begin{latin}
\url{http://en.wikipedia.org/wiki/Bit_operation}
\end{latin}

همچنین مقاله‌ای پیرامون عملگرهای سه‌گانه (\en{ternary operations}) ارائه کرده است:

\begin{latin}
\url{http://en.wikipedia.org/wiki/Ternary_operation}
\end{latin}

افزون بر این، شرحی بر فرمول محاسبه اقساط وام که در اسکریپت \cmd{loan-calc} از آن استفاده کرده‌ایم، در دسترس است:

\begin{latin}
\url{http://en.wikipedia.org/wiki/Amortization_calculator}
\end{latin}